\markdownRendererDocumentBegin
\markdownRendererWarning{The `hybrid` option has been soft-deprecated.}{The `hybrid` option has been soft-deprecated.}{Consider using one of the following better options for mixing TeX and markdown: `contentBlocks`, `rawAttribute`, `texComments`, `texMathDollars`, `texMathSingleBackslash`, and `texMathDoubleBackslash`. For more information, see the user manual at <https://witiko.github.io/markdown/>.}{Consider using one of the following better options for mixing TeX and markdown: `contentBlocks`, `rawAttribute`, `texComments`, `texMathDollars`, `texMathSingleBackslash`, and `texMathDoubleBackslash`. For more information, see the user manual at <https://witiko.github.io/markdown/>.}\markdownRendererSectionBegin
\markdownRendererHeadingOne{Meet in the Middle}\markdownRendererInterblockSeparator
{}El algoritmo "Meet in the Middle" es una técnica de optimización que reduce la complejidad computacional de problemas de búsqueda exhaustiva mediante la partición del espacio de soluciones en dos mitades independientes. Su fundamento radica en el principio divide y vencerás: en lugar de explorar todas las $2^n$ combinaciones posibles en un problema, se divide el conjunto de n elementos en dos subconjuntos de tamaño n/2, generando todas las configuraciones posibles para cada mitad ($2^(n/2)$ cada una), almacenando los resultados de la primera mitad en una estructura de datos eficiente (típicamente una tabla hash o árbol de búsqueda), y luego explorando la segunda mitad buscando complementos que satisfagan las condiciones del problema.\markdownRendererInterblockSeparator
{}\markdownRendererInputFencedCode{./_markdown_main/aecf0f243d83a828ecf7d3fa75891c1f.verbatim}{cpp}{cpp}
\markdownRendererSectionEnd \markdownRendererDocumentEnd